\addhighlights{yes}
\renewcommand{\addhighlights}{

\noindent\textbf{\boldmath{What are the main findings?}}

\begin{enumerate}[label=$\bullet$]

\item FP2MP, a blackboard-based multi-agent LLM architecture, solves ill-defined urban planning problems via iterative uncertainty reduction, achieving superior metacognition (5.00 $\pm$ 0.00 vs. 2.26-3.51 for baselines) and uncertainty handling (4.74 $\pm$ 0.45 vs. 2.46-3.83) across 60 problem instances. The FP2MP‑Eval protocol demonstrates high inter‑rater reliability and sensitivity to model capability differences, enabling multi-dimensional assessment of problem structuring, solution quality, and reflective reasoning without ground‑truth answers.

\end{enumerate}

\noindent\textbf{\boldmath{What is the implication of the main findings?}}

\begin{enumerate}[label=$\bullet$]

\item The proposed framework introduces FP2MP, a blackboard-based LLM-MAS architecture that operationalises evidence-based ill-defined urban planning problem solving by combining tool-augmented agents with a shared three-level memory and critic-driven closed-loop control. Proposed FP2MP-Eval evaluation framework provides a multi-dimensional LLM-as-a-judge evaluation protocol for ill-defined tasks, enabling reproducible assessment of problem structuring, solution quality, and reflective reasoning without ground truth. The open‑source implementation (core architecture, evaluation, and baselines) offers a reusable foundation for developing and benchmarking LLM‑MAS on ill-defined problems beyond urban planning.

\end{enumerate}
}